% Generated draft: requires mathematical and source review.
\documentclass[11pt]{article}
\usepackage[margin=25mm]{geometry}
\usepackage{fontspec}
\setmainfont{Latin Modern Roman}
\usepackage{amsmath,amssymb,mathtools,unicode-math}
\setmathfont{Latin Modern Math}
\usepackage{enumitem,needspace,xurl,hyperref}
\setlength{\parindent}{0pt}
\setlength{\parskip}{5pt}
\setlength{\emergencystretch}{3em}
\newcommand{\N}{\mathbb{N}}
\newcommand{\Z}{\mathbb{Z}}
\newcommand{\Q}{\mathbb{Q}}
\newcommand{\R}{\mathbb{R}}
\newcommand{\C}{\mathbb{C}}
\newcommand{\F}{\mathbb{F}}
\newcommand{\A}{\mathbb{A}}
\newcommand{\PP}{\mathbb{P}}
\newcommand{\E}{\mathbb{E}}
\newcommand{\eps}{\varepsilon}
\newcommand{\dd}{\,\mathrm{d}}
\begin{document}
\section*{UnsolvedMath ID: 30000766}
\textbf{Source problem code:} catalogue entry
\subsection{Definitions and mathematical statement}
Let $D\subset\PP^2_{\Q}$ be a normal-crossing divisor of degree at least four, let $X=\PP^2_{\Q}\setminus D$, fix an affine embedding $\iota:X\hookrightarrow\A^n_{\Q}$, and let $S$ be a finite set of rational primes. Put
\[
 \Z[S^{-1}]=\{a/b\in\Q:\text{every prime dividing }b\text{ lies in }S\}
\]
and define the integral points with respect to this fixed embedding by
\[
 X(\Z[S^{-1}];\iota)
 =\{x\in X(\Q):\iota(x)\in\Z[S^{-1}]^n\}.
\]
The conjecture is that, for every allowed $D$, embedding, and finite $S$,
\[
 \overline{X(\Z[S^{-1}];\iota)}^{\mathrm{Zar}}\ne X.
\]
Equivalently the integral points lie in a proper Zariski-closed subset, allowed to depend on $D$, $\iota$, and $S$. The conclusion is non-density, not finiteness of the point set or a uniform exceptional subset for all $S$ and embeddings.

The usual reduced normal-crossings convention means that, after passage to an algebraic closure, the divisor is locally a smooth branch or two smooth transverse branches in surface coordinates. The source does not explicitly state reducedness, multiplicity conventions, or simple versus ordinary normal crossings. This matters: degree of a divisor with multiplicities and degree of its support need not agree, while the complement depends only on the support. The intended reduced/support-degree hypothesis must be checked rather than silently counted with arbitrary multiplicities.

No splitting of the components over $\Q$, line-arrangement assumption, or restriction on $|S|$ is imposed. The supplied assessment places the question in higher-dimensional nonabelian-Chabauty motivation but calls it a finiteness question; that wording is not substituted for the explicit weaker non-Zariski-density target.
\subsection{Short English statement}
For a normal-crossing plane-divisor complement of degree at least four, must the rational points integral outside a fixed finite set of primes in a chosen affine embedding lie in a proper algebraic subset? The claim is non-density, with reducedness and divisor-degree conventions needing confirmation.
\subsection{Sources}
\begin{itemize}
\item Local UnsolvedMath registry, permanent ID 30000766, code catalogue entry, ``Integral-Point Non-Density on Plane-Divisor Complements.'' Proposed year 2007; proposer unspecified.
\item \emph{Arithmetic and Differential Galois Groups} (2008), Oberwolfach Reports, pp.\ 1443--1520, DOI 10.4171/owr/2007/26; \url{https://doi.org/10.4171/owr/2007/26}. The displayed report year 2008 differs from the DOI and proposed year 2007. This is the sole supplied literature reference.
\item Statement status corrected\_verified; the supplied verification says the official report was checked directly, the complete non-density question restored, and this is the canonical extraction of a repeated passage.
\item Literature triage checked 2026-08-21: open. It describes higher-dimensional nonabelian-Chabauty motivation and a general unresolved finiteness question. The exact statement here is integral-point non-density, and the relation to that broader wording needs review. Created 2026-08-25; updated 2026-09-16. Current status and source claims require independent verification.
\item Registry curation by Ulam AI / UnsolvedMath Contributors, CC BY 4.0; the underlying report retains its supplied provenance. Unreviewed statement-only draft.
\end{itemize}
\end{document}
